% =============================================================
% Capitolo: User Profiling
% Progetto: MCP_ECOM — Agente e-commerce basato su MCP
% =============================================================

\chapter{User Profiling}
\label{chap:user-profiling}

% ─────────────────────────────────────────────────────────────
\section{Introduzione}
\label{sec:profiling-intro}
% ─────────────────────────────────────────────────────────────

Il sistema di user profiling in MCP\_ECOM ha l'obiettivo di rendere
l'agente \emph{progressivamente più personale} ad ogni interazione,
senza richiedere configurazioni esplicite da parte dell'utente. Il
principio guida è la \textbf{coesistenza}: le preferenze apprese
automaticamente non sovrascrivono mai le impostazioni scelte
consapevolmente dall'utente, ma le arricchiscono.

Il profilo si articola su due assi ortogonali:

\begin{description}
    \item[Preferenze esplicite] Impostate dall'utente tramite il tool MCP
        \texttt{update\_user\_preferences}: brand favoriti, fascia di
        prezzo, istruzioni personalizzate, tono conversazionale, tema UI.

    \item[Preferenze auto-apprese] Derivate automaticamente dalle
        interazioni di ricerca: affinità di categoria, budget contestuale,
        condizione preferita (nuovo / usato / ricondizionato), profilo
        comportamentale.
\end{description}

Tutto il profilo viene persistito su \textbf{PostgreSQL}, sincronizzato
nella \textbf{memoria a lungo termine} Redis ad ogni richiesta, ed esposto
all'agente tramite una \textbf{risorsa MCP} dedicata.


% ─────────────────────────────────────────────────────────────
\section{Schema del Profilo Utente}
\label{sec:profiling-schema}
% ─────────────────────────────────────────────────────────────

Il modello SQLAlchemy \texttt{User} (\texttt{app/models/user.py})
definisce due gruppi di colonne distinti.

\subsection{Preferenze Esplicite}
\label{subsec:explicit-prefs}

\begin{lstlisting}[language=Python,caption={Colonne di preferenza esplicita nel modello User (app/models/user.py)},label=lst:user-model-explicit]
favorite_brands     = Column(String(255), nullable=True)
price_preference    = Column(String(50),  nullable=True)
custom_instructions = Column(String(1000),nullable=True)
conversation_tone   = Column(String(50),  default="neutral")
theme               = Column(String(20),  default="light")
\end{lstlisting}

\subsection{Campi Auto-Appresi (Retrocompatibili)}
\label{subsec:auto-learned-cols}

I quattro campi di apprendimento automatico sono stati aggiunti in una
fase successiva del progetto e sono dichiarati con \texttt{nullable=True}
e \texttt{server\_default=None}, garantendo piena retrocompatibilità con
gli utenti pre-esistenti:

\begin{lstlisting}[language=Python,caption={Colonne auto-apprese nel modello User (app/models/user.py)},label=lst:user-model-auto]
# JSON: {"Cellulari & smartphone": 12, "Scarpe": 5}
category_affinities = Column(String(2000), nullable=True,
                              server_default=None)

# Formato counter: "new:5,used:2" oppure valore piatto legacy "new"
condition_preference = Column(String(20), nullable=True,
                               server_default=None)

# Tier: "browser" | "researcher" | "power_buyer"
interaction_depth   = Column(String(20), nullable=True,
                              server_default=None)

# JSON: {"auto_cat:scarpe": 150.0, "auto_brand:nike": 200.0}
contextual_budgets  = Column(String(2000), nullable=True)
\end{lstlisting}

La scelta di serializzare strutture complesse come stringhe JSON all'interno
di colonne \texttt{VARCHAR} — anziché usare colonne \texttt{JSONB} di
PostgreSQL — è stata adottata per massimizzare la portabilità verso
qualunque backend SQL.


% ─────────────────────────────────────────────────────────────
\section{Pipeline di Auto-Apprendimento}
\label{sec:auto-learning}
% ─────────────────────────────────────────────────────────────

Il punto di ingresso del sistema di apprendimento è la funzione
\texttt{update\_user\_profile()} (\texttt{app/services/user\_profiling.py}),
chiamata dalla search pipeline al termine di ogni ricerca andata a buon
fine. La funzione è \textbf{non-committing}: modifica l'oggetto ORM ma
lascia al chiamante la responsabilità del \texttt{db.commit()}, evitando
conflitti di transazione.

\begin{lstlisting}[language=Python,caption={Chiamata a update\_user\_profile() dalla search pipeline (app/services/search\_pipeline.py)},label=lst:auto-learn-trigger]
changed = update_user_profile(
    user=user,
    parsed=parsed,                   # Output del parser semantico
    db=db,
    dominant_category_name=_dominant_cat_name,
    items=items[:10],
    action="search"
)
if changed:
    db.commit()
\end{lstlisting}

La funzione orchestra cinque sotto-aggiornamenti indipendenti, ciascuno
isolato in un blocco \texttt{try/except} per garantire che il fallimento
di uno non comprometta gli altri.

\begin{figure}[htbp]
    \centering
    \fbox{\parbox{0.88\linewidth}{\centering\vspace{2.2cm}
        \textit{[Diagramma: query utente → parser → search pipeline →
        update\_user\_profile() → 5 sotto-funzioni → PostgreSQL commit]}\\
        \vspace{2.2cm}}}
    \caption{Pipeline di auto-apprendimento del profilo utente.}
    \label{fig:profiling-pipeline}
\end{figure}

\subsection{Brand Preference}
\label{subsec:brand-learning}

I brand vengono appresi esclusivamente dai brand \emph{validati} dal
parser semantico (campo \texttt{parsed["brands"]}), non da estrazioni
euristiche sul testo grezzo. I nuovi brand vengono preposti alla lista
esistente (LIFO), con deduplicazione case-insensitive e cap a 10 elementi:

\begin{lstlisting}[language=Python,caption={Aggiornamento brand preference (app/services/user\_profiling.py)},label=lst:brand-update]
# Prepend nuovi brand, deduplica, cap a 10
merged = brands + [b for b in current_list
                   if b.lower() not in {nb.lower() for nb in brands}]
user.favorite_brands = ",".join(merged[:10])
\end{lstlisting}

Le scelte manuali dell'utente vengono sempre preservate: il meccanismo
\emph{aggiunge} nuovi brand in testa, ma non ne rimuove mai di esistenti.

\subsection{Global Price Preference — EMA}
\label{subsec:price-ema}

La fascia di prezzo globale viene aggiornata tramite una
\textbf{Exponential Moving Average} (EMA) con fattore di apprendimento
$\alpha = 0.20$:

\begin{equation}
    p_{\text{new}} = \lfloor p_{\text{old}} \times 0.80 + p_{\text{signal}} \times 0.20 \rfloor
    \label{eq:ema-price}
\end{equation}

Questo meccanismo smussante evita che un singolo acquisto anomalo (es.
una ricerca spot per un prodotto di lusso) distorca eccessivamente il
profilo di budget abituale.

\subsection{Category Affinities}
\label{subsec:category-affinities}

Ogni ricerca incrementa il contatore della categoria dominante
(derivata dai risultati eBay). Le affinità sono memorizzate come
dizionario JSON ordinato per conteggio decrescente, con cap a 20
categorie:

\begin{lstlisting}[language=Python,caption={Aggiornamento category affinities (app/services/user\_profiling.py)},label=lst:cat-affinities]
affinities[cat] = int(affinities.get(cat, 0)) + 1

# Mantieni solo le top 20 categorie
sorted_cats = sorted(affinities.items(),
                     key=lambda x: x[1], reverse=True)
affinities = dict(sorted_cats[:20])

user.category_affinities = json.dumps(affinities)
# Esempio: {"Cellulari & smartphone": 12, "Scarpe": 5, "Laptop": 3}
\end{lstlisting}

\subsection{Contextual Budgets}
\label{subsec:contextual-budgets}

Questo è il meccanismo più granulare: viene mantenuto un budget EMA
separato per categoria e per brand, usando chiavi con prefisso
\texttt{auto\_} per distinguerle dalle voci manuali (che non vengono
mai sovrascritte):

\begin{lstlisting}[language=Python,caption={Aggiornamento contextual budgets con prefisso auto\_ (app/services/user\_profiling.py)},label=lst:contextual-budgets]
# Budget per categoria
cat_key = f"auto_cat:{dominant_category_name.lower()}"
if cat_key in budgets:
    # EMA: 80% vecchio + 20% nuovo segnale
    budgets[cat_key] = round(old * 0.80 + price_signal * 0.20, 2)
else:
    budgets[cat_key] = round(price_signal, 2)

# Budget per brand
brand_key = f"auto_brand:{brand.lower()}"
# ... stesso aggiornamento EMA ...

# Risultato in DB:
# {"auto_cat:scarpe": 175.0, "auto_brand:nike": 200.0,
#  "manual_premium_brands": 500.0}  <- voci manuali intatte
\end{lstlisting}

\subsection{Condition Preference — Majority Counter}
\label{subsec:condition-counter}

La condizione preferita (nuovo / usato / ricondizionato) viene determinata
tramite un \textbf{contatore a maggioranza}: si commette una preferenza
solo dopo almeno 3 segnali, di cui il candidato dominante deve
rappresentare oltre il 50\% del totale:

\begin{lstlisting}[language=Python,caption={Logica majority counter per condition preference (app/services/user\_profiling.py)},label=lst:condition-counter]
# Formato storage: "new:5,used:2"
counter[condition] = counter.get(condition, 0) + 1
total = sum(counter.values())

# Soglia: >= 3 segnali, share > 50%
if total >= 3:
    top = max(counter, key=lambda k: counter[k])
    if counter[top] / total >= 0.5:
        dominant = top   # <- preferenza confermata

user.condition_preference = ",".join(f"{k}:{v}"
                                      for k, v in sorted(counter.items()))
\end{lstlisting}

Il formato counter viene anche usato per l'override manuale: impostare
\texttt{"new:99"} tramite il tool MCP forza la dominanza di
\emph{nuovo} indipendentemente dai segnali precedenti.

\subsection{Interaction Depth — State Machine}
\label{subsec:interaction-depth}

Il profilo comportamentale dell'utente avanza attraverso tre tier con
una state machine deterministica. Il tier viene alzato in base
all'azione eseguita:

\begin{table}[htbp]
\centering
\begin{tabular}{llll}
\hline
\textbf{Tier corrente} & \textbf{Azione} & \textbf{Nuovo tier}
    & \textbf{Score azione} \\
\hline
\texttt{browser}    & \texttt{detail / compare / seller} & \texttt{researcher}   & 3 / 3 / 2 \\
\texttt{browser}    & \texttt{contact}                   & \texttt{power\_buyer} & 5 \\
\texttt{researcher} & \texttt{contact}                   & \texttt{power\_buyer} & 5 \\
qualsiasi           & \texttt{search}                    & invariato             & 1 \\
\hline
\end{tabular}
\caption{State machine dell'interaction depth.}
\label{tab:interaction-depth}
\end{table}

I tre tier hanno un significato semantico preciso che l'agente usa
nella fase di final answer:

\begin{description}
    \item[\texttt{browser}] Utente che esplora, prevalentemente ricerche.
    \item[\texttt{researcher}] Utente approfondito: analizza venditori,
        legge specifiche tecniche, confronta modelli.
    \item[\texttt{power\_buyer}] Acquirente esperto: contatta venditori,
        tratta prezzi, ha un'alta frequenza di interazione.
\end{description}


% ─────────────────────────────────────────────────────────────
\section{Gestione Manuale delle Preferenze}
\label{sec:manual-prefs}
% ─────────────────────────────────────────────────────────────

L'utente può aggiornare esplicitamente qualunque campo del profilo
tramite il tool MCP \texttt{update\_user\_preferences}
(\texttt{app/mcp/tools/profile.py}). Il tool accetta i seguenti
parametri opzionali:

\begin{center}
\begin{tabular}{lp{8cm}}
\hline
\textbf{Parametro} & \textbf{Effetto} \\
\hline
\texttt{favorite\_brands}           & Sovrascrive i brand preferiti \\
\texttt{price\_preference}          & Imposta il budget di riferimento \\
\texttt{custom\_instructions}       & Regole di sistema personalizzate \\
\texttt{conversation\_tone}         & Stile conversazionale dell'agente \\
\texttt{theme}                      & Tema UI (\texttt{light} / \texttt{dark}) \\
\texttt{condition\_preference}      & Override condizione con \texttt{"value:99"} \\
\texttt{reset\_category\_affinities} & Cancella le affinità apprese \\
\hline
\end{tabular}
\end{center}

Dopo ogni aggiornamento manuale, il tool notifica il frontend via
\texttt{notify\_resource\_updated("profile://\{user\_id\}")} in modo
che l'interfaccia rifletta immediatamente il nuovo stato.


% ─────────────────────────────────────────────────────────────
\section{Esposizione del Profilo all'Agente}
\label{sec:profile-exposure}
% ─────────────────────────────────────────────────────────────

Il profilo raggiunge l'LLM attraverso due canali distinti che si
completano a vicenda.

\subsection{Canale 1 — Risorsa MCP \texttt{profile://}}
\label{subsec:mcp-resource-channel}

Il profilo viene servito come risorsa MCP dall'endpoint
\texttt{profile://\{session\_id\}} (\texttt{app/mcp/resources.py}).
La risorsa costruisce il JSON di risposta fondendo i campi espliciti
con quelli auto-appresi elaborati da \texttt{build\_profile\_context()}:

\begin{lstlisting}[language=Python,caption={Endpoint risorsa MCP profile (app/mcp/resources.py)},label=lst:mcp-resource]
@mcp.resource("profile://{session_id}")
async def get_user_profile(session_id: str) -> str:
    profile_data = {
        "user_id": user.id,
        "favorite_brands": user.favorite_brands,
        "price_preference": user.price_preference,
        "conversation_tone": user.conversation_tone,
        "language": user.language,
        ...
    }
    # Arricchisce con i campi auto-appresi
    profile_ctx = build_profile_context(user)
    profile_data["condition_preference"] = profile_ctx["condition_preference"]
    profile_data["interaction_depth"]    = profile_ctx["interaction_depth"]
    profile_data["category_affinities"]  = profile_ctx["category_affinities"]
    profile_data["contextual_budgets"]   = profile_ctx["contextual_budgets"]
    return _safe_json({"status": "ok", "profile": profile_data})
\end{lstlisting}

La funzione \texttt{build\_profile\_context()} effettua due
post-elaborazioni rilevanti prima di esporre il dato:

\begin{itemize}
    \item \textbf{Category affinities}: restituisce solo le top~3 categorie
        ordinate per conteggio, per non sovraccaricare il prompt.
    \item \textbf{Contextual budgets}: rimuove il prefisso \texttt{auto\_}
        dalle chiavi, rende il JSON leggibile dall'LLM e mantiene
        separate le voci manuali.
\end{itemize}

\subsection{Canale 2 — Iniezione nelle Custom Instructions}
\label{subsec:injection-channel}

All'avvio di ogni sessione, l'agente (\texttt{app/agent/ebay\_agent.py})
recupera la risorsa MCP e ne traduce i campi principali in istruzioni
testuali, aggiungendole alle \texttt{custom\_instructions} dell'utente
con il tag \texttt{[MCP PROFILE]}:

\begin{lstlisting}[language=Python,caption={Fetch e iniezione del profilo MCP nell'agente (app/agent/ebay\_agent.py)},label=lst:agent-profile-inject]
raw_profile = await self.mcp_client.read_resource_async(
    f"profile://{user_id}"
)
prof_data = json.loads(raw_profile).get("profile", {})

addons = []
if prof_data.get("language"):
    addons.append(
        f"Rispondi sempre all'utente nella lingua "
        f"'{prof_data['language']}'"
    )
if prof_data.get("favorite_brands"):
    addons.append(
        f"L'utente preferisce questi brand: {prof_data['favorite_brands']}"
    )
if prof_data.get("price_preference"):
    addons.append(
        f"Note sul budget dell'utente: {prof_data['price_preference']}"
    )

extra = " ".join(addons)
# Aggiunge al custom_instructions esistente
setattr(self.user, "custom_instructions",
        current + " | [MCP PROFILE] " + extra)
\end{lstlisting}

Oltre al profilo utente, l'agente recupera anche il prompt MCP
\texttt{shopping\_expert\_prompt} e lo prepone alle custom instructions,
fornendo all'LLM l'identità base di \emph{assistente e-commerce
professionale} prima ancora delle preferenze dell'utente.

\subsection{Canale 3 — Long-Term Memory nello Scratchpad}
\label{subsec:ltm-channel}

La memoria a lungo termine (\texttt{app/agent/memory.py}) sincronizza
i campi del profilo DB nella struttura \texttt{long\_term\_memory}
dello scratchpad all'inizio di ogni richiesta. Questo terzo canale
è quello consumato dalla funzione \texttt{build\_final\_answer\_prompt()}
per personalizzare la risposta finale:

\begin{lstlisting}[language=Python,caption={Struttura long\_term\_memory nello scratchpad},label=lst:ltm-structure]
long_term_memory = {
    "user_preferences": {
        "db_favorite_brands":      "Nike,Adidas",
        "db_price_preference":     "175",
        "db_condition_preference": "new",      # dominant se >= 3 segnali
        "db_interaction_depth":    "researcher",
        "validated_brands":        ["Nike"],   # parser-validated (sessione)
        "category_history":        ["Scarpe", "Abbigliamento"],
        "favorite_sellers":        ["monclick"],
        "recent_brand_hints":      ["Nike"]    # legacy token hints
    },
    "user_behaviour": {
        "search_count": 12,
        "interaction_counts": {
            "search": 10, "detail": 2,
            "compare": 1, "seller": 1, "contact": 0
        }
    },
    "previous_searches": ["nike air max", "scarpe running"]
}
\end{lstlisting}


% ─────────────────────────────────────────────────────────────
\section{Wishlist e Tracking dei Prezzi}
\label{sec:wishlist}
% ─────────────────────────────────────────────────────────────

Il sistema di wishlist (\texttt{app/models/wishlist.py},
\texttt{app/mcp/tools/wishlist\_tool.py}) è strettamente integrato con
il profilo utente: consente di salvare prodotti eBay di interesse e di
ricevere notifiche di variazione del prezzo.

\begin{lstlisting}[language=Python,caption={Campi principali del modello WishlistItem},label=lst:wishlist-model]
id, user_id (FK -> users.id),
ebay_id, title, price, previous_price,   # price tracking
currency, condition, image_url, url,
seller_name,
added_at, last_checked_at
\end{lstlisting}

Il campo \texttt{previous\_price} permette di rilevare riduzioni di
prezzo tra un controllo e il successivo. Il tool MCP
\texttt{manage\_wishlist} supporta le operazioni \texttt{add},
\texttt{remove} e \texttt{list}, con sincronizzazione asincrona in
background verso la Watchlist ufficiale di eBay.


% ─────────────────────────────────────────────────────────────
\section{Flusso Dati Complessivo}
\label{sec:profiling-dataflow}
% ─────────────────────────────────────────────────────────────

\begin{figure}[htbp]
    \centering
    \fbox{\parbox{0.9\linewidth}{\centering\vspace{3cm}
        \textit{[Diagramma: query → parser → search pipeline →
        update\_user\_profile() → PostgreSQL →
        hydrate\_request\_state() → Redis (LTM) →
        MCP resource fetch → custom\_instructions →
        planner + final answer]}\\
        \vspace{3cm}}}
    \caption{Flusso dati del sistema di user profiling.}
    \label{fig:profiling-dataflow}
\end{figure}

Il profilo percorre il seguente ciclo ad ogni richiesta:

\begin{enumerate}
    \item \textbf{Apprendimento} (fine ricerca): \texttt{update\_user\_profile()}
        aggiorna i campi ORM e il chiamante esegue \texttt{db.commit()}.

    \item \textbf{Idratazione} (inizio richiesta successiva):
        \texttt{hydrate\_request\_state()} sincronizza i campi DB nella
        struttura \texttt{long\_term\_memory} in Redis (TTL 30 giorni).

    \item \textbf{Fetch MCP} (avvio sessione agente): l'agente legge
        \texttt{profile://\{user\_id\}}, costruisce le istruzioni
        personalizzate e le inietta come prefisso nelle
        \texttt{custom\_instructions}.

    \item \textbf{Planning personalizzato}: il planner riceve le
        \texttt{custom\_instructions} arricchite e le usa come strato
        a priorità massima nel prompt.

    \item \textbf{Final answer personalizzata}: \texttt{build\_final\_answer\_prompt()}
        legge la \texttt{long\_term\_memory} e inietta nel prompt
        brand, budget, condizione, categorie, profilo comportamentale
        e ricerche precedenti.
\end{enumerate}


% ─────────────────────────────────────────────────────────────
\section{Riepilogo dei Meccanismi}
\label{sec:profiling-summary}
% ─────────────────────────────────────────────────────────────

\begin{table}[htbp]
\centering
\small
\begin{tabular}{p{3.2cm} p{2.5cm} p{3cm} p{4cm}}
\hline
\textbf{Campo} & \textbf{Tipo} & \textbf{Algoritmo} & \textbf{Note} \\
\hline
\texttt{favorite\_brands}
    & Esplicito + Auto
    & Prepend + dedup (cap 10)
    & Solo brand validati dal parser \\
\hline
\texttt{price\_preference}
    & Esplicito + Auto
    & EMA $\alpha=0.20$
    & Budget globale smussato \\
\hline
\texttt{category\_affinities}
    & Auto
    & Hit counter (cap 20)
    & Top 3 esposte al prompt \\
\hline
\texttt{contextual\_budgets}
    & Esplicito + Auto
    & EMA per chiave (prefisso \texttt{auto\_})
    & Manuale e auto coesistono \\
\hline
\texttt{condition\_preference}
    & Esplicito + Auto
    & Majority counter
    & Commit dopo $\geq3$ segnali, $>50\%$ \\
\hline
\texttt{interaction\_depth}
    & Auto
    & State machine (3 tier)
    & browser → researcher → power\_buyer \\
\hline
\texttt{custom\_instructions}
    & Esplicito
    & Concatenazione
    & Priorità assoluta nel prompt \\
\hline
\texttt{conversation\_tone}
    & Esplicito
    & Lookup table
    & Iniettato nel final answer \\
\hline
\end{tabular}
\caption{Riepilogo dei meccanismi di aggiornamento del profilo utente.}
\label{tab:profiling-mechanisms}
\end{table}

Il sistema garantisce tre proprietà architetturali fondamentali:
\textbf{retrocompatibilità} (colonne nullable, degradazione graceful),
\textbf{coesistenza} (chiave \texttt{auto\_} per non sovrascrivere le
scelte manuali) e \textbf{isolamento transazionale} (nessun commit
dentro \texttt{update\_user\_profile()}, delegato al chiamante).
